\chapter{Cough}

\section*{Definition}
Reflex action clearing airways of mucus or irritants; classified acute, subacute, chronic by duration.

\section*{When to See a Doctor}
See your doctor if:
\begin{itemize}
  \item Symptoms lasting more than 3 weeks (persistent)
\end{itemize}

\section*{How It Develops}
Productive vs dry, timing, triggers help differentiate etiations; upper airway cough syndrome most common cause of chronic cough.

\section*{Causes}
\begin{itemize}
  \item Viral URI
  \item Asthma/COPD
  \item GERD
  \item Postnasal drip
  \item Smoking/smoke irritation
\end{itemize}

\section*{Related Symptoms}
\begin{itemize}
  \item Sore throat
  \item Congestion
  \item Wheezing
  \item Chest tightness
  \item Phlegm production
\end{itemize}

\section*{Medical Evaluation}
\begin{itemize}
  \item History focuses on onset, duration, character of symptoms (productive vs dry cough, nasal discharge color), fever pattern, and exposure history.
  \item Physical examination includes vital signs, lung auscultation, nasal inspection, and oropharyngeal examination.
  \item Targeted testing may include rapid antigen/ PCR for respiratory viruses, chest X-ray, pulse oximetry, and sputum culture.
  \item Allergy testing or pulmonary function tests are reserved for chronic or recurrent cases.
\end{itemize}

\section*{Treatment Options}
\begin{itemize}
  \item Treatment depends on the cause; uncomplicated viral cough is usually managed with fluids, rest, and symptom relief.
  \item Antibiotics are not routinely useful for viral cough and should be prescribed only when a bacterial cause is suspected or confirmed.
  \item Asthma, COPD, reflux, postnasal-drip, and chronic cough require cause-specific treatment and assessment when persistent.
\end{itemize}

\section*{Self-Care and Home Management}
\begin{itemize}
  \item Maintain adequate hydration with water, warm tea, or broth to thin secretions.
  \item Use a humidifier or steam inhalation to soothe irritated airways and loosen mucus.
  \item Rest and avoid exposure to smoke, dust, and other respiratory irritants.
  \item Over-the-counter remedies such as saline nasal spray, cough drops, or honey (for adults) may provide symptomatic relief.
  \item Monitor for warning signs including high fever, difficulty breathing, chest pain, or worsening symptoms.
\end{itemize}

\section*{References}
\begin{thebibliography}{99}
\bibitem{cough:cough1} Irwin RS, Baumann MH, Bolser DC, et al. ``Diagnosis and management of cough: ACCP evidence-based clinical practice guidelines.'' Chest. 2006;129(1 Suppl):1S-23S.
\bibitem{cough:cough2} Gibson PG, Vertigan AE. ``Management of chronic cough.'' BMJ. 2015;351:h5590.
\bibitem{cough:cough3} Morice AH, Millqvist E, Bieksiene K, et al. ``ERS guidelines on the diagnosis and treatment of chronic cough in adults.'' Eur Respir J. 2020;55(1):1901136.
\bibitem{cough:cough4} Irwin RS, French CT, Chang AB, et al. ``Classification of cough as a symptom in adults and children.'' Chest. 2018;153(1):196-209.
\bibitem{cough:cough5} Kamei RK, Weisman SJ. ``Pertussis (whooping cough).'' In: Nelson Textbook of Pediatrics. 21st ed. Elsevier; 2020.
\bibitem{cough:cough6} Braman SS. ``Postinfectious cough: ACCP evidence-based clinical practice guidelines.'' Chest. 2006;129(1 Suppl):138S-146S.
\end{thebibliography}
